% TOP_PROBLEM: {"id": 1451, "title": "Eilenberg-Ganea Conjecture", "category_id": 4, "category": {"id": 4, "name": "algebra", "display_name": "Algebra", "description": "Group theory, ring theory, field theory, and algebraic structures.", "slug": "algebra", "order_index": 4, "created_at": "2026-07-31T15:26:25.671Z"}, "status": "open"}
% FIRST_POSTED: 2026-10-01
% ATTEMPT_STATUS: unresolved
% Imported from: unsolvedmath_additional_500.tex; UnsolvedMath ID 1451
% !TeX program = lualatex
% Compile twice: lualatex 1451.tex
% Self-contained: no external bibliography, images, or \input files.
% Numeric section labels and visible numbers are original UnsolvedMath IDs.
\documentclass[11pt,a4paper]{article}
\usepackage[margin=24mm,headheight=14pt]{geometry}
\usepackage{fontspec}
\setmainfont{Latin Modern Roman}
\setsansfont{Latin Modern Sans}
\setmonofont{Latin Modern Mono}
\usepackage{amsmath,amssymb,mathtools}
\usepackage{unicode-math}
\setmathfont{Latin Modern Math}
\usepackage{microtype}
\usepackage{enumitem,xurl,needspace}
\usepackage[dvipsnames]{xcolor}
\usepackage{hyperref}
\hypersetup{unicode=true,colorlinks=true,linkcolor=MidnightBlue,urlcolor=MidnightBlue,
 pdftitle={UnsolvedMath Problem 1451},pdfauthor={Source-annotated compilation},bookmarksnumbered=true}
\usepackage{bookmark,tocloft,titlesec,fancyhdr}
\setlength{\cftsecnumwidth}{4.2em}
\cftsetpnumwidth{2.5em}
\cftsetrmarg{4em}
\titleformat{\section}{\Large\bfseries\raggedright}{\thesection}{0.7em}{}
\pagestyle{fancy}\fancyhf{}
\fancyhead[L]{\small\sffamily Additional UnsolvedMath statements}
\fancyhead[R]{\small Original catalogue IDs}
\fancyfoot[C]{\thepage}
\setlength{\parindent}{0pt}
\setlength{\parskip}{5pt plus 1pt}
\setlength{\emergencystretch}{3em}
\setcounter{tocdepth}{1}\setcounter{secnumdepth}{1}
\setlist[itemize]{leftmargin=3em,itemsep=4pt,topsep=4pt,parsep=0pt}
\allowdisplaybreaks
\newcommand{\N}{\mathbb{N}}
\newcommand{\Z}{\mathbb{Z}}
\newcommand{\Q}{\mathbb{Q}}
\newcommand{\R}{\mathbb{R}}
\newcommand{\C}{\mathbb{C}}
\newcommand{\F}{\mathbb{F}}
\newcommand{\A}{\mathbb{A}}
\newcommand{\PP}{\mathbb{P}}
\newcommand{\E}{\mathbb{E}}
\newcommand{\eps}{\varepsilon}
\newcommand{\dd}{\,\mathrm d}
\newcommand{\sref}[2]{\hyperlink{source:#1:#2}{[S#2]}}
\newif\ifshowcatalogue
\showcataloguetrue
\newcommand{\cataloguescope}[1]{\ifshowcatalogue
 \paragraph{Statement recorded in the supplied source.}
 #1\fi}
\begin{document}
% UnsolvedMath ID 1451; source catalogue code ALG-006

\setcounter{section}{1450}

\section{The Eilenberg--Ganea Dimension-Two Conjecture}\label{1451}

{\small\sffamily UnsolvedMath ID 1451 · Algebra}

\subsection{Definitions and mathematical statement}

\paragraph{Shared notation.}Unless a section says otherwise, $\N=\{1,2,\ldots\}$,
$\Z$ is the ring of integers, $\Q,\R,\C$ are the rational, real, and complex
fields, $[N]=\{1,\ldots,\lfloor N\rfloor\}$, and $\log$ is the natural
logarithm. For real functions of a parameter tending to infinity,
$f\ll g$ and $f=O(g)$ mean $|f|\leq Cg$ eventually for some fixed $C$;
$f\gg g$ reverses this comparison; $f=o(g)$ means $f/g\to0$;
$f\sim g$ means $f/g\to1$; and $f\asymp g$ means both order bounds.
A subscript on $O$ or $\ll$ specifies allowed constant dependence.
The expression $n^{o(1)}$ in an upper bound means growth smaller than
$n^\epsilon$ for each fixed $\epsilon>0$. Local definitions govern any
exception, finite-field convention, or use of the same symbol for another
object. An asymptotic claim is not automatically an effective algorithm.



The integral cohomological dimension $\operatorname{cd}_{\Z}G$ is the projective dimension of the trivial $\Z G$-module $\Z$. A $K(G,1)$ is a connected CW complex with fundamental group $G$ and contractible universal cover. The geometric dimension is the least dimension of such a complex. Ask whether $\operatorname{cd}_{\Z}G=2$ forces geometric dimension two. \sref{1451}{1} \sref{1451}{2}

\paragraph{Statement recorded in the supplied source.}
Does every group with cohomological dimension 2 have a 2-dimensional Eilenberg-MacLane space K(G,1)? \sref{1451}{1}

\paragraph{Residual task reported by the archive.}

The embedded review (2026-08-17) records: Determine the geometric dimension of the Bestvina--Brady candidate groups, equivalently overcoming the associated Whitehead-asphericity obstruction, or prove a general two-dimensional realization theorem. \sref{1451}{1}

\subsection{Short English statement}

If a group has integral cohomological dimension two, must it admit a genuinely two-dimensional classifying space?

\subsection{Research attempt: an explicit dimension-two resolution and product-of-trees models}
\textbf{Status.} We prove the conclusion for $\mathbb Z^2$ and for products of two nontrivial finitely generated free groups. We do not realize arbitrary projective resolutions geometrically. The primary Bestvina--Brady paper, inspected through [R1], studies finiteness properties of subgroups of right-angled Artin groups; it does not permit identifying all cohomological and geometric conditions without further work.

\subsubsection{1. An exact algebraic calculation for the torus group}
Let $R=\mathbb Z[x^{\pm1},y^{\pm1}]$, the group ring of $\mathbb Z^2$. Consider
\[0\longrightarrow R\xrightarrow{d_2}R^2\xrightarrow{d_1}R
\xrightarrow{\epsilon}\mathbb Z\longrightarrow0,
\quad d_2(c)=((1-y)c,(x-1)c),\quad
d_1(a,b)=(x-1)a+(y-1)b.\tag{1}\]
Here $\epsilon$ evaluates both variables at 1. Its kernel is $(x-1,y-1)$, so exactness at the last $R$ follows. If $d_1(a,b)=0$, reducing modulo $x-1$ gives $(y-1)\overline b=0$ in the integral domain $\mathbb Z[y^{\pm1}]$. Thus $b=(x-1)c$; cancellation in $R$ then gives $a=(1-y)c$. This proves exactness at $R^2$. The first map is injective because $x-1$ is not a zero divisor. Hence (1) is a free resolution of length two.

Applying $\operatorname{Hom}_R(-,\mathbb Z)$ with the trivial action makes both differentials zero, since $x-1$ and $y-1$ act by zero. It follows that $H^2(\mathbb Z^2;\mathbb Z)=\mathbb Z$, proving the projective dimension cannot be smaller than two. The torus is independently a two-dimensional $K(\mathbb Z^2,1)$: its universal cover is $\mathbb R^2$, with the translation action. Thus both dimensions are exactly two, and the resolution is the cellular chain calculation of an actual space.

\subsubsection{2. A larger family with genuine geometric models}
Let $F_r,F_s$ be free groups with $r,s\geq1$, and take graphs $R_r,R_s$ consisting of bouquets of $r,s$ circles. Their universal covers are trees. The square complex $R_r\times R_s$ has fundamental group $F_r\times F_s$ and universal cover the product of those trees, which contracts coordinatewise. It is therefore a two-dimensional classifying space.

The elementary cellular chain complex of a bouquet has zero differential and first homology $\mathbb Z^r$. The product chain complex gives $H^2(R_r\times R_s;\mathbb Z)=\mathbb Z^{rs}\ne0$. Equivalently this is the torsion-free Kunneth calculation. Thus cohomological dimension is at least two and the geometric model bounds it above by two. Both dimensions are exactly two throughout this family.

\subsubsection{3. The realization step that cannot be skipped}
The hypothesis in the question supplies a length-two \emph{projective} resolution of a module. A CW universal cover supplies a based \emph{free} chain complex, with boundary maps arising from actual attaching maps, together with contractibility. These are additional requirements. Arbitrary ring matrices need not already be boundary maps of a chosen presentation complex, and vanishing ordinary homology of a proposed universal cover is not sufficient until simple connectivity and the required homotopical conditions are established.

The exactness proof of (1) succeeds because it comes with a torus model; it is not a template for replacing all projectives by two-dimensional cells. Likewise a subgroup of a group acting on a contractible complex may act on a complex of a larger dimension than its cohomological dimension. Restricting the action does not reduce the dimension automatically.

\subsubsection{4. Verification and residual question}
The proof audit checks each kernel and image in (1), including Laurent rather than polynomial variables, and the nonvanishing cohomology lower bound. No numerical test can verify arbitrary group resolutions. The product examples also show why a two-dimensional space can have many degree-two cohomology classes without requiring higher-dimensional cells. What remains is a general realization theorem that turns the given projective-dimension condition into a contractible two-dimensional universal cover, or a rigorously established group for which every classifying space requires dimension at least three.

\subsection{Sources}

{\small

\begin{itemize}

\item[\hypertarget{source:1451:1}{S1}] ULAM.AI, \emph{UnsolvedMath}, supplied \texttt{problems.json}, record 1451 (ALG-006), statement and background. The embedded literature-review date is 2026-08-17; that is the archive’s review date, not a fresh certification. Dataset landing page: \url{https://huggingface.co/datasets/ulamai/UnsolvedMath}.

\item[\hypertarget{source:1451:2}{S2}] Mladen Bestvina and Noel Brady, Morse theory and finiteness properties of groups, Inventiones Mathematicae 129 (1997), 445-470. \url{https://doi.org/10.1007/s002220050168}. Bibliographic information imported from the supplied record.

\item[\hypertarget{source:1451:3}{S3}] James Howie, Bestvina--Brady groups and the plus construction, Mathematical Proceedings of the Cambridge Philosophical Society 127 (1999), 487-493. \url{https://doi.org/10.1017/S0305004199003928}. Bibliographic information imported from the supplied record.

\item[\hypertarget{source:1451:4}{S4}] A. Adem and I. Hambleton, Minimal Euler characteristics for even-dimensional manifolds with finite fundamental group, Forum of Mathematics, Sigma 11 (2023), e24. \url{https://doi.org/10.1017/fms.2023.18}. Bibliographic information imported from the supplied record.

\item[R1] M. Bestvina and N. Brady, Morse theory and finiteness properties of groups, Inventiones Mathematicae 129 (1997), primary paper.
\url{https://people.math.osu.edu/davis.12/courses/8800-20/Bestvina-Brady.pdf}.
\end{itemize}

}

\label{end:1451}
\end{document}
